\section{Case Study: A Zone Monitor}
\label{sec:evaluation}

\texttt{ZoneMonitor} is a B machine implemented and imported into \barrel
for this case study. It manages readings from redundant sensors and
maintains authorization for zones whose readings satisfy a configured
policy. Partial-map application, cardinality, and extrema arise from the
policy and its operations. The development exercises the translation of
these expressions, automatic WD discharge, reusable relational lemmas,
and individual proofs of invariant preservation. All 146 retained
obligations are proved and published in Lean without admissions.

\subsection{State and safety policy}

Let $S$ and $Z$ denote the finite nonempty carrier sets \texttt{SENSORS}
and \texttt{ZONES}. A fixed total map $\mathit{zone}\in S\to Z$ assigns
each sensor to a zone, and $q\in Z\to\mathbb{N}_{>0}$ specifies the
required number of available sensors. Measurements lie in a fixed nonempty
integer interval $D$. The state contains
\[
  r\in S\pfun D,\qquad
  \ell,u\in Z\to D,\qquad P\subseteq Z.
\]
Here $r$ contains the available readings, $\ell$ and $u$ give the configured
bounds, and $P$ contains the zones currently authorized. The B variables
are named \texttt{reading}, \texttt{lower}, \texttt{upper}, and
\texttt{permit}, respectively. Define
\[
  A_z=\operatorname{dom}(r)\cap\mathit{zone}^{-1}[\{z\}],
  \qquad V_z=r[A_z].
\]
The quorum counts sensors in $A_z$, not distinct values in $V_z$: two sensors
reporting the same value still constitute two available sensors. The invariant
requires $\ell(z)\le u(z)$ for every $z\in Z$, together with
\begin{equation}
\label{eq:zone-policy}
\forall z\in P,\quad
\begin{aligned}
  &A_z\ne\varnothing\ \land\ q(z)\le\operatorname{card}(A_z)\ \land{}\\
  &\ell(z)\le\min(V_z)\ \land\ \max(V_z)\le u(z).
\end{aligned}
\end{equation}
Non-emptiness precedes the extrema in the invariant. Finiteness of $S$
makes $A_z$ finite, and $A_z\subseteq\operatorname{dom}(r)$ makes its image
nonempty whenever $A_z$ is nonempty. No assumption requires every zone to
have enough installed sensors: a zone with too few can never be authorized.

Initialization sets $r=P=\varnothing$ and assigns initial bound maps
$\ell_0,u_0\in Z\to D$, supplied as constants satisfying
$\ell_0(z)\le u_0(z)$. The policy concerns recorded measurements; it does
not assume that sensors report the true physical values. The resulting
property is authorization consistency within the monitor, not a
certification of the monitored plant.

\subsection{Operations and the batch-update proof}

Table~\ref{tab:zone-operations} summarizes the eight operations. Reading
updates and configuration changes revoke affected authorizations. This
allows readings to change freely while preventing a permit from retaining
a justification based on stale values.

\begin{table}[tb]
\caption{Operations of the implemented \texttt{ZoneMonitor} machine.}
\label{tab:zone-operations}
\small
\begin{tabularx}{\textwidth}{@{}lX@{}}
\toprule
Operation & Guard or state change \\
\midrule
\texttt{Record} & Override $r$ at sensor $s$ and remove $\mathit{zone}(s)$ from $P$. \\
\texttt{RecordBatch} & Override $r$ with $b\in S\pfun D$ and remove $\mathit{zone}[\operatorname{dom}(b)]$ from $P$. \\
\texttt{Withdraw} & Require $s\in\operatorname{dom}(r)$; remove its reading and revoke its zone. \\
\texttt{Configure} & Set the bounds for $z$, requiring the new lower bound not to exceed the upper; revoke $z$. \\
\texttt{Authorize} & Require the conditions in (\ref{eq:zone-policy}) for $z\in Z$; add $z$ to $P$. \\
\texttt{Revoke} & Remove $z\in Z$ from $P$. \\
\texttt{ReadSensor} & Require $s\in\operatorname{dom}(r)$; return $r(s)$. \\
\texttt{ZoneSummary} & Require $z\in Z$ and $A_z\ne\varnothing$; return its cardinality and the minimum and maximum of $V_z$. \\
\bottomrule
\end{tabularx}
\end{table}

The batch operation is written directly in B:
\begin{lstlisting}[language={},morekeywords={PRE,THEN,END}]
RecordBatch(batch) =
  PRE batch : SENSORS +-> VALUE
  THEN
    reading := reading <+ batch ||
    permit := permit - zone[dom(batch)]
  END
\end{lstlisting}
Here \texttt{+->} is the ASCII form of B's partial-function-space arrow
$\pfun$, \texttt{<+} is relational override, and \texttt{||} combines the
two assignments in parallel. Thus both right-hand sides refer to the
old state. The operation updates any subset of sensor readings in one
step, and preserves permits only for zones outside the image of the
updated sensors.

Write $r'=r\mathbin{\oplus}b$ for the new reading relation and
$T=\mathit{zone}[\operatorname{dom}(b)]$. For $z\in P\setminus T$, no
sensor assigned to $z$ belongs to $\operatorname{dom}(b)$. Relational
override leaves both the presence and the value of every reading in that
zone unchanged. The Lean development proves the corresponding frame
lemmas \texttt{available\_overload} and \texttt{readings\_overload}:
\[
  z\notin T\quad\Longrightarrow\quad
  A'_z=A_z\ \land\ r'[A'_z]=r[A_z].
\]
These lemmas are stated for arbitrary relations over arbitrary types;
their proofs use set extensionality, the domain of override, and the
definition of relational image. They do not depend on sensor bounds or
on the authorization invariant. Non-emptiness, quorum, and the two
extremum bounds can then be transported separately to the updated state.

For example, the following is the actual \barrel command proving the
lower-bound preservation obligation:
\begin{lstlisting}[mathescape=true]
obligation Operation_RecordBatch_15 of ZoneMonitor by
  intros
  expose_names
  obtain $\langle$_hne, _hcard, hmin, _hmax$\rangle$ := h_15 zz h_17.1
  simpa only [readings_overload h_17.2] using hmin
\end{lstlisting}
The generated hypothesis \texttt{h\_15} is the old authorization policy;
\texttt{h\_17.1} states that the selected zone was permitted, and
\texttt{h\_17.2} that it is unaffected by the batch. The proof extracts
the old lower bound and rewrites the updated image using the frame lemma.
The maximum and quorum obligations use the same argument. Single-sensor
updates reduce to the singleton case, while withdrawal uses analogous
lemmas for domain subtraction. Configuration changes additionally require
lemmas describing application after a singleton override.

The WD prerequisite of this lower-bound obligation is a separate named
goal. Its proof first rewrites the updated image and then applies the
finite-nonempty-image minimum lemma, using non-emptiness from the old
policy. This order exposes the distinction between establishing that the
minimum exists and proving its lower bound. Both arguments reuse the same
mathematical fact about unaffected zones, but they establish different
propositions.

\subsection{Source WD and generated WD obligations}

The artifact imports a POG generated by Atelier~B with its \texttt{-w}
option. It contains 41 principal obligations and 39 source WD obligations.
The latter include definedness of constants and the invariant, operation
preconditions, and output expressions. This is significant for
\texttt{ReadSensor} and \texttt{ZoneSummary}: these read-only operations
preserve the state, but their application and extrema still require WD
proofs. Their five source WD goals are included in the development.

These source obligations differ from the conditions created by \barrel
when it translates a partial expression. Source WD obligations are B
predicates emitted by Atelier~B. Generated WD obligations are proof
arguments required by the Lean operator definitions. Importing the full
POG therefore proves both families. It also avoids counting the absence
of an invariant-preservation goal for an output-only operation as evidence
that its output expressions are defined. The case uses
\texttt{import pog ZoneMonitor from "specs/zonemonitor"}; the ordinary
machine-import path currently invokes the generator without \texttt{-w}.

Encoding the 80 source goals creates 754 requests for WD evidence.
Sharing retains 66 distinct generated WD goals and reuses 688 earlier
conditions. The repeated requests arise from the same typing hypotheses,
policy clauses, and updated-state expressions across generated goals.
They do not correspond to 754 distinct mathematical facts. The retained
66 conditions, together with the 80 source goals, give 146 obligations
in the \barrel development.

\subsection{Automation and remaining proof work}

The case supplies proved helper lemmas and a bounded WD tactic branch
through \barrel's extension mechanism. The helpers derive finiteness of
available sensors, membership in a map's source, and non-emptiness from
an authorization policy. They then invoke the library's application,
cardinality, and finite-image extrema lemmas. The branch recognizes the
four corresponding WD predicates and uses bounded searches for their
premises. This is the measured configuration; the results are not a claim
about unextended stock automation.

On import, automation discharges 51 of the 66 distinct generated WD goals
($77.3\%$) and 20 of the 80 source goals. The latter comprise 10 principal
and 10 source WD goals. The remaining 75 obligations receive individual
proof commands: 31 principal, 29 source WD, and 15 generated WD goals.
The 29 source WD commands explicitly invoke a reusable
\texttt{zone\_source\_wd} tactic. They are counted as supplied proof
commands, not as successes of import-time automatic discharge. The
remaining generated WD proofs use the frame lemmas and authorization
case distinctions described above.

\begin{table}[tb]
\caption{Obligations by source. Parentheses give import-time automatic
proofs. Generated WD goals are assigned to the source that first introduces
them; later uses contribute requests but no new distinct goal.}
\label{tab:zone-evidence}
\small
\begin{tabular}{@{}lrrrr@{}}
\toprule
Source & Principal & Source WD & WD requests & Unique WD \\
\midrule
Properties & 0 & 4 (4) & 0 & 0 \\
Initialization & 5 (4) & 0 & 16 & 8 (8) \\
Invariant & 0 & 15 (3) & 76 & 4 (4) \\
\texttt{Record} & 6 (1) & 2 (0) & 91 & 15 (13) \\
\texttt{RecordBatch} & 6 (1) & 0 & 66 & 6 (4) \\
\texttt{Withdraw} & 6 (1) & 2 (0) & 91 & 7 (5) \\
\texttt{Configure} & 8 (1) & 0 & 88 & 8 (6) \\
\texttt{Authorize} & 5 (1) & 11 (3) & 220 & 12 (7) \\
\texttt{Revoke} & 5 (1) & 0 & 56 & 6 (4) \\
\texttt{ReadSensor} & 0 & 1 (0) & 10 & 0 \\
\texttt{ZoneSummary} & 0 & 4 (0) & 40 & 0 \\
\midrule
Total & 41 (10) & 39 (10) & 754 & 66 (51) \\
\bottomrule
\end{tabular}
\end{table}

Table~\ref{tab:zone-evidence} reports the counts by source. In particular,
the read-only operations contribute 50 generated requests, all reusing
conditions introduced earlier. Their zero in the final column therefore
indicates sharing, not an absence of definedness reasoning. The evidence
export retains the generated declaration dependencies so that the shared
conditions and their consumers can be inspected.

The completed development uses Lean~4.31.0, with a budget of 20,000
heartbeats per import-time proof attempt and 200,000 for subsequent
proof commands. Running \texttt{qed ZoneMonitor} publishes all 146
declarations. A subsequent audit checks that there are no skipped,
pending, or admitted obligations and examines the transitive axioms of
every published obligation. The only axioms found are Lean's standard
\texttt{propext}, \texttt{Classical.choice}, and \texttt{Quot.sound}.
The artifact retains the B machine, full POG, helper lemmas, individual
proof commands, JSON evidence, tabulated counts, run log, and source
hashes. Thus the counts describe an executed proof development.

The automatic fraction therefore measures routine proof work removed at
import time under a documented extension. The residual proofs identify
where the application policy matters: preserving readings outside an
update, recovering non-emptiness from a permit, and distinguishing an
existing permit from a newly authorized zone. This single machine tests
those interactions without a refinement chain or a separate industrial
benchmark. Industrial inputs can supply additional evidence about import
scale; their size is independent of the proof-completion result for this
case. Kernel checking establishes the translated propositions under the
embedding, with the translation trust boundary stated in
Section~\ref{sec:workflow}.
